\chapter{Manual de Configuración del Plugin}
\label{Appendix:Configuracion}

Este apéndice describe todos los parámetros configurables del plugin desde el panel de propiedades de OBS Studio. Para abrirlo, hacer clic derecho sobre la fuente de vídeo $\rightarrow$ \textit{Filtros} $\rightarrow$ seleccionar el filtro del plugin.

El parámetro principal es el desplegable \textbf{Modo de renderizado}, que determina qué controles son visibles. Los cinco modos disponibles son: \textbf{3D}, \textbf{AR}, \textbf{Countdown (reloj)}, \textbf{Scoreboard (texto)} y \textbf{Team Info}.

% -------------------------------------------------------
\section*{Modo 3D}

Renderiza un modelo 3D en una posición fija en pantalla, sin detección de marcadores.

\begin{itemize}
  \item \textbf{Ruta del Modelo 3D} — fichero del modelo a cargar (formatos: \texttt{.obj}, \texttt{.fbx}, \texttt{.dae}, \texttt{.gltf}).
  \item \textbf{Ruta de la Textura} — imagen de textura a aplicar sobre el modelo (\texttt{.png}, \texttt{.jpg}, \texttt{.bmp}, \texttt{.tga}). Opcional si el modelo lleva textura embebida.
  \item \textbf{Escala} — factor de escala uniforme del modelo (rango $0.1$--$1000$).
  \item \textbf{Posicion X / Y / Z} — posición del modelo en el espacio de pantalla (rango $-3000$ a $3000$, paso $10$).
  \item \textbf{Rotacion X / Y / Z (Grados)} — orientación del modelo en cada eje (rango $-360°$ a $360°$).
\end{itemize}

% -------------------------------------------------------
\section*{Modo AR}

Detecta un marcador ArUco en el vídeo y ancla el modelo 3D sobre él en tiempo real. La posición manual no está disponible en este modo; la posición la determina el marcador detectado.

\subsection*{Modelo 3D}

\begin{itemize}
  \item \textbf{Ruta del Modelo 3D} — mismo comportamiento que en modo 3D.
  \item \textbf{Ruta de la Textura} — mismo comportamiento que en modo 3D.
  \item \textbf{Escala} — factor de escala del modelo.
\end{itemize}

\subsection*{Detección ArUco}

\begin{itemize}
  \item \textbf{ID del Marker} — identificador numérico del marcador ArUco que el plugin debe detectar (rango $0$--$100$). Debe coincidir exactamente con el ID del marcador impreso.
  \item \textbf{Tamano del Marker} — lado real del marcador en metros (rango $0.1$--$10.0$). Un valor incorrecto afecta directamente a la estimación de posición y escala.
  \item \textbf{Diccionario de Marker} — familia de marcadores ArUco utilizada. Opciones disponibles: \texttt{Original}, \texttt{4x4 (100)}, \texttt{5x5 (100)}, \texttt{6x6 (100)}, \texttt{7x7 (100)}. Debe coincidir con el diccionario usado al generar el marcador.
  \item \textbf{Archivo de Calibracion} — ruta al fichero \texttt{.yml} con los parámetros intrínsecos de la cámara generado por la herramienta de calibración incluida con el plugin. Si se deja vacío se usan valores por defecto que pueden producir desviaciones en el posicionamiento.
\end{itemize}

\subsection*{Ajuste fino (offsets)}

Permiten corregir pequeños desvíos visuales sin necesidad de recalibrar la cámara:

\begin{itemize}
  \item \textbf{AR Offset Posicion X / Y / Z} — desplazamiento adicional del modelo respecto al centro del marcador (rango $-1000$ a $1000$).
  \item \textbf{AR Offset Rotacion X / Y / Z} — corrección de orientación sobre los ejes del marcador (rango $-360°$ a $360°$).
\end{itemize}

% -------------------------------------------------------
\section*{Modo Countdown (reloj)}

Muestra el modelo 3D cargado actuando como reloj analógico de cuenta atrás. Puede posicionarse de forma fija en pantalla o anclarse a un marcador ArUco.

\subsection*{Modelo y posición}

\begin{itemize}
  \item \textbf{Ruta del Modelo 3D}, \textbf{Ruta de la Textura}, \textbf{Escala} — igual que en los modos anteriores.
  \item \textbf{Reloj en Marker AR} — casilla que determina el modo de posicionamiento. Si está \textbf{desactivada}, aparecen los controles de posición manual (Posicion X/Y/Z y Rotacion X/Y/Z). Si está \textbf{activada}, esos controles se ocultan y aparecen en su lugar los controles de detección ArUco (ID del Marker, Tamaño, Diccionario, Archivo de Calibración y Offsets), funcionando igual que en el modo AR.
\end{itemize}

\subsection*{Configuración del reloj}

\begin{itemize}
  \item \textbf{Horas / Minutos / Segundos} — duración total de la cuenta atrás.
  \item \textbf{Manecillas Reloj} — modo de animación de las mallas del modelo:
    \begin{itemize}
      \item \textbf{H/M/S (3)} — tres manecillas independientes (hora, minuto, segundo). La manecilla de segundos avanza $6°$ por segundo, la de minutos $6°$ por minuto completado, y la de horas de forma proporcional al rango configurado.
      \item \textbf{Total (1)} — una única manecilla que barre $360°$ a lo largo de toda la duración configurada, representando el progreso total del concurso.
    \end{itemize}
  \item \textbf{Cuenta Atras a Mover} — casilla para iniciar o pausar la cuenta atrás.
  \item \textbf{Reiniciar Reloj} — casilla para restablecer el tiempo al valor inicial configurado.
\end{itemize}

\subsection*{Sincronización con DOMjudge}

Visible en este modo y en los modos Scoreboard y Team Info. Ver la sección \textit{Conexión con DOMjudge} más adelante.

% -------------------------------------------------------
\section*{Modo Scoreboard (texto)}

Muestra la tabla de clasificación del concurso descargada de DOMjudge. No requiere marcador físico para posicionarse.

\subsection*{Posición}

\begin{itemize}
  \item \textbf{Centrar en Pantalla} — si está activado, el scoreboard se centra automáticamente en la escena y los offsets manuales se ignoran.
  \item \textbf{Offset Manual X / Y} — posición del scoreboard cuando el centrado automático está desactivado (rango $0$--$3000$).
\end{itemize}

\subsection*{Apariencia del texto}

\begin{itemize}
  \item \textbf{Tamano de Fuente} — tamaño en puntos del texto (rango $5$--$150$).
  \item \textbf{Color del Texto} — color de las letras.
  \item \textbf{Color del Borde} — color del contorno de cada carácter.
  \item \textbf{Grosor del Borde} — anchura del contorno en píxeles (rango $0$--$20$).
  \item \textbf{Max caracteres equipo} — número máximo de caracteres del nombre del equipo antes de truncarlo (rango $5$--$40$). Solo disponible en modo Scoreboard.
  \item \textbf{Bandas por fila} — activa el sombreado alternado entre filas pares e impares. Solo disponible en modo Scoreboard.
  \item \textbf{Opacidad bandas (\%)} — transparencia de las bandas de fila (rango $0$--$80$). Solo disponible en modo Scoreboard.
\end{itemize}

\subsection*{Fondo del overlay}

\begin{itemize}
  \item \textbf{Activar fondo del texto} — muestra un rectángulo de fondo semitransparente detrás del texto.
  \item \textbf{Color de fondo} — color del rectángulo.
  \item \textbf{Opacidad fondo (\%)} — transparencia del fondo (rango $0$--$100$).
  \item \textbf{Margen interno (px)} — espacio entre el borde del rectángulo y el texto (rango $0$--$80$).
  \item \textbf{Radio esquinas (px)} — redondeo de las esquinas del rectángulo (rango $0$--$80$).
\end{itemize}

\subsection*{Sombra del overlay}

\begin{itemize}
  \item \textbf{Activar sombra} — dibuja una sombra paralela detrás del fondo.
  \item \textbf{Opacidad sombra (\%)} — transparencia de la sombra (rango $0$--$100$).
  \item \textbf{Sombra offset X / Y (px)} — desplazamiento horizontal y vertical de la sombra (rango $-50$ a $50$).
  \item \textbf{Sombra suavidad (capas)} — número de pasadas de difuminado; valores mayores producen sombras más suaves (rango $1$--$8$).
\end{itemize}

\subsection*{Suavizado de posición AR}

Estos controles están presentes en los modos Scoreboard y Team Info cuando el overlay se ancla a un marcador:

\begin{itemize}
  \item \textbf{Suavizado AR (overlay)} — activa el filtro exponencial de suavizado de posición para reducir el temblor del overlay.
  \item \textbf{Suavizado AR (alpha)} — factor del filtro. Valores próximos a $1.0$ siguen el marcador con mayor fidelidad; valores próximos a $0.0$ producen mayor inercia y movimiento más suave (rango $0.01$--$1.0$).
\end{itemize}

\subsection*{Sincronización con DOMjudge}

Ver sección siguiente.

% -------------------------------------------------------
\section*{Modo Team Info}

Detecta cualquier marcador registrado en el fichero JSON de mapeo y muestra los datos del equipo asociado. Comparte los controles de apariencia (texto, fondo, sombra, suavizado) con el modo Scoreboard, con las siguientes diferencias: no tiene posición manual ni opción de centrado, no tiene límite de caracteres por equipo ni bandas por fila, ya que muestra únicamente un equipo cada vez.

\begin{itemize}
  \item \textbf{Diccionario de Marker} — familia ArUco utilizada para detectar todos los marcadores del fichero JSON.
  \item \textbf{JSON de equipos (ArUco$\rightarrow$Team)} — ruta al fichero \texttt{.json} con el mapeo de marcadores a equipos. El formato esperado es:
\begin{verbatim}
[
  { "marker_id": 1, "team_id": "42" },
  { "marker_id": 2, "team_id": "12" }
]
\end{verbatim}
  donde \texttt{marker\_id} es el ID del marcador ArUco colocado en la mesa del equipo y \texttt{team\_id} el identificador del equipo en DOMjudge. El fichero admite hasta 16 entradas.
  \item \textbf{Recargar JSON} — botón que recarga el fichero de mapeo en tiempo de ejecución sin necesidad de reiniciar el filtro ni la escena.
\end{itemize}

% -------------------------------------------------------
\section*{Conexión con DOMjudge}

Los modos \textbf{Countdown}, \textbf{Scoreboard} y \textbf{Team Info} pueden sincronizarse con DOMjudge. Los parámetros son idénticos en los tres modos:

\begin{itemize}
  \item \textbf{Sincronizar con DOMjudge} (\textit{Sincronizacion Web} en modo Countdown) — activa la descarga periódica de datos desde la API.
  \item \textbf{URL API (DOMjudge)} — URL base de la instancia de DOMjudge, por ejemplo:
\begin{verbatim}
https://miservidor.domjudge.org/api/v4
\end{verbatim}
    No incluir la barra final. Este es el campo que varía en cada evento y debe actualizarse antes de cada retransmisión.
  \item \textbf{Contest ID} — identificador numérico del concurso activo en DOMjudge, visible en la URL de la interfaz web del concurso.
  \item \textbf{Usuario API} — nombre de usuario con permisos de lectura en la API de DOMjudge.
  \item \textbf{Contrasena API} — contraseña del usuario anterior (campo oculto).
  \item \textbf{Probar Conexion} — botón que realiza una solicitud de prueba a la API. El resultado se muestra en el log de OBS (\textit{Ayuda} $\rightarrow$ \textit{Registro}).
  \item \textbf{Intervalo Sinc (seg)} / \textbf{Intervalo API (seg)} — frecuencia de actualización automática en segundos (rango $1$--$300$). Un valor de $30$ segundos es adecuado para la mayoría de concursos.
\end{itemize}